% ============================================================================
% Chapter 3 — Dialectal / Nabati Poetry (الشعر العامي / الشعر النبطي)
% ============================================================================
\chapter{Dialectal and Nabati Poetry}
\label{ch:nabati}

Nabati poetry (\ar{الشعر النبطي}) — also called ``the people's poetry'' (\ar{شعر الشعب}) — is the great vernacular literary tradition of the Arabian Peninsula. Composed in spoken Arabian dialects rather than Standard Arabic, it has maintained an unbroken oral tradition of at least six centuries while developing its own rigorous prosodic system parallel to (and partially derived from) the Khalilian system of classical Arabic.

% ============================================================================
\section{Origins and Historical Development}
\label{sec:nabati_history}

The earliest documented Nabati poetry dates to the 14th century CE (referenced by Ibn Khaldun in his \textit{Muqaddimah}, 1377). The tradition flowered from the 16th century onward as Arabian Peninsula dialects diverged increasingly from literary Arabic. Unlike classical poetry, which was composed for literate court audiences, Nabati poetry was composed by and for tribal Bedouin communities — often by unschooled poets who had absorbed prosodic intuitions from the oral environment rather than from formal study.

The name ``Nabati'' has contested etymology: some trace it to the ancient Nabataean Arabs; others argue it derives from \ar{نَبَط} (to draw water from a well) — suggesting verse that ``draws out'' emotion. The term is now used without etymological weight simply to mean ``Arabian Peninsula colloquial verse.''

\subsection{Relationship to Khalilian Prosody}

The Nabati prosodic system is \textbf{not a degenerate version of Khalilian prosody}. It is a parallel development within the same phonological family. Key differences:
\begin{itemize}
  \item Nabati meters are realised through \textit{colloquial phonology}: glottal stops replace hamzas; vowel shifts (e.g., \ar{قلب} ``heart'' → \ar{غلب} in Najdi) affect syllable lengths.
  \item The \textbf{dual-rhyme system} (a distinct rhyme for each hemistich) has no equivalent in classical poetry.
  \item Nabati meters permit more syllabic flexibility in the hashw (middle feet) than classical meters.
  \item Tashkeel (full diacritics) is less systematically applied to Nabati text — the phonological realisation is implicit in the dialect.
\end{itemize}

% ============================================================================
\section{The Nabati Prosodic System}
\label{sec:nabati_prosody}

Nabati scholars have identified \textbf{twelve to fourteen} recognised meters (\ar{بحور نبطية}), divided into \textbf{primary} (\ar{رئيسية}) and \textbf{innovated} (\ar{مبتكرة}) categories.

\subsection{Primary Nabati Meters (\ar{البحور الرئيسية})}

\begin{table}[H]
\centering
\caption{Primary Nabati meters}
\label{tab:nabati_primary_meters}
\renewcommand{\arraystretch}{1.6}
\begin{tabularx}{\textwidth}{llXlX}
\toprule
\textbf{Name} & \textbf{Arabic} & \textbf{Classical Analogue} & \textbf{Rhyme Type} & \textbf{Use} \\
\midrule
Al-Hilali & \ar{الهلالي} & Taweel variant & Single (classical) & All themes; oldest Nabati meter \\
Al-Sakhri & \ar{الصخري} & Wafir & Single & Narrative, tribal history \\
Al-Hida' & \ar{الحداء} & Rajaz & Single & Camel-driving, martial \\
Al-Rajad & \ar{الرجد} & Mutadarak & Single & Light verse, songs \\
Al-Mashhub & \ar{المسحوب} & Hilali-derived (shorter) & \textbf{Dual} & The dominant Nabati meter (\textgreater60\%) \\
Al-Hijini & \ar{الهجيني} & Sari'/Baseet mixed & Single & Camel-riding, journey verse \\
\bottomrule
\end{tabularx}
\end{table}

\subsection{Innovated Nabati Meters (\ar{البحور المبتكرة})}

\begin{table}[H]
\centering
\caption{Innovated Nabati meters}
\label{tab:nabati_invented_meters}
\renewcommand{\arraystretch}{1.6}
\begin{tabularx}{\textwidth}{llXX}
\toprule
\textbf{Name} & \textbf{Arabic} & \textbf{Foot Pattern (approximate)} & \textbf{Associated Tradition} \\
\midrule
Al-Mashhub & \ar{المسحوب} & \ar{مستفعلن مستفعلن فاعلاتن} (each hemistich) & General / Samri \\
Bahr al-Samiri & \ar{بحر السامري} & Linked to Mashhub; nocturnal gathering meter & Samri performance \\
Al-Zuhayri & \ar{الزهيري} & \ar{مستفعلن فاعلن مستفعلن فاعل} & Dense rhyming verse \\
Bahr al-Funun & \ar{بحر الفنون} & \ar{فاعلاتن فاعلاتن فاعلن} & Innovated by Ibn La'bun \\
Al-Murabbi' & \ar{المربع / المربوع} & 4-line strophic; AAA-B rhyme & Muwashshahat-like forms \\
Al-Jinas & \ar{الجناس} & Based on wordplay patterns & Occasional, humorous \\
\bottomrule
\end{tabularx}
\end{table}

% ============================================================================
\section{The Dual-Rhyme System (\ar{نظام القافية المزدوجة})}
\label{sec:dual_rhyme}

The most distinctive feature of the Mashhub meter — and of the dominant strand of Nabati poetry — is the \textbf{dual-rhyme system}: each bayt has \textit{two distinct rhyme letters}, one for the end of the first hemistich (\ar{الصدر}) and a \textit{different} one for the end of the second hemistich (\ar{العجز}).

\begin{rulebox}{Rules of the Nabati Dual-Rhyme System}
\begin{enumerate}
  \item The \textbf{sadr rhyme} (\ar{قافية الصدر}) applies consistently to all first hemistichs throughout the poem.
  \item The \textbf{ajuz rhyme} (\ar{قافية العجز}) applies consistently to all second hemistichs throughout the poem.
  \item The two rhymes \textbf{must never converge} — if the sadr rhyme ends with \ar{ن} (nun), the ajuz rhyme must not also end with \ar{ن}. Convergence is considered a prosodic error in the Mashhub tradition.
  \item The dual-rhyme system is \textbf{unique to Nabati poetry}; importing it into classical Fusha verse is an error.
  \item Both rhymes must be internally consistent throughout the poem (no switching rawiy mid-poem).
\end{enumerate}
\end{rulebox}

\begin{verseframe}
\centering\ar{يا صاحبي اللي يبالي بوالفي \quad\quad والقلب مني يتمنى لقاكم}\\[4pt]
\ar{قافية الصدر: فاء (ـِي) \quad قافية العجز: كاف (ـكم)}\\[2pt]
\small\textit{Example demonstrating dual rhyme: sadr rhyme on \ar{ف} / ajuz rhyme on \ar{ك}}
\end{verseframe}

% ============================================================================
\section{The Hilali and Mashhub Meters in Detail}
\label{sec:hilali_mashhub}

\subsection{Al-Hilali (\ar{الهلالي})}

The oldest documented Nabati meter, named for the Banu Hilal tribe. It is the ancestor of the Mashhub and uses a single (classical-style) rhyme.

\begin{table}[H]
\centering
\caption{Al-Hilali meter structure}
\renewcommand{\arraystretch}{1.4}
\begin{tabularx}{\textwidth}{lX}
\toprule
\textbf{Hemistich pattern (approximate)} & \ar{مستفعلن مستفعلن فاعلاتن} (long form) \\
\textbf{Rhyme type} & Single rhyme (classical-style) \\
\textbf{Association} & Oldest Nabati poems; heroic/historical verse; tribal epics \\
\textbf{Derivation} & Related to classical Taweel but realised through Najdi/Hijazi phonology \\
\textbf{Modern use} & Increasingly rare; considered archaic by contemporary Nabati poets \\
\bottomrule
\end{tabularx}
\end{table}

\subsection{Al-Mashhub (\ar{المسحوب})}

The dominant Nabati meter — over 60\% of all Nabati poetry is in the Mashhub. Its name derives from \ar{سَحَبَ} (to draw/pull) — referring to the way a singer ``draws'' the note along with the bow of the rababa (bowed string instrument).

\begin{table}[H]
\centering
\caption{Al-Mashhub meter structure}
\label{tab:mashhub}
\renewcommand{\arraystretch}{1.4}
\begin{tabularx}{\textwidth}{lX}
\toprule
\textbf{Hemistich pattern} & \ar{مُسْتَفْعِلُنْ مُسْتَفْعِلُنْ فَاعِلَاتُنْ} \\
\textbf{Derivation} & Shortened from al-Hilali (one syllable shorter); emerged in the 12th century AH \\
\textbf{Rhyme type} & \textbf{Dual rhyme} (sadr rhyme ≠ ajuz rhyme) — mandatory \\
\textbf{Associated tradition} & The Samri nocturnal performance; general Nabati composition \\
\textbf{Credit for innovation} & Muhsin al-Hazani (\ar{محسن الحزاني}) credited with formalising the Mashhub and dual-rhyme system in Samri tradition \\
\textbf{Musical character} & Designed for singing with melodic rababa accompaniment; the shortened final foot (\ar{فاعلاتن}) creates the characteristic ``pull'' effect \\
\bottomrule
\end{tabularx}
\end{table}

\begin{rulebox}{AI Generation Rules for Al-Mashhub}
\begin{enumerate}
  \item Choose two \textit{distinct} rhyme letters before generating — one for sadr, one for ajuz.
  \item Apply the dual-rhyme constraint to every verse without exception.
  \item The hemistich pattern \ar{مستفعلن مستفعلن فاعلاتن} allows syllabic flexibility in the middle feet similar to classical Rajaz.
  \item The final foot (\ar{فاعلاتن}) must be respected — it carries the rhyme and the musical ``resolution.''
  \item Language must be in the appropriate Arabian Peninsula dialect — classical Fusha is a mismatch.
\end{enumerate}
\end{rulebox}

% ============================================================================
\section{Nabati Poetic Purposes (\ar{أغراض الشعر النبطي})}
\label{sec:nabati_aghard}

Nabati poetry shares most classical أغراض but with a different emphasis and some unique categories:

\begin{table}[H]
\centering
\caption{Nabati poetic purposes}
\label{tab:nabati_aghard}
\renewcommand{\arraystretch}{1.5}
\begin{tabularx}{\textwidth}{llX}
\toprule
\textbf{Purpose} & \textbf{Arabic} & \textbf{Notes and Distinctives} \\
\midrule
Tribal Praise / Chivalry & \ar{شهامة وفروسية} & The Nabati equivalent of madih; praises tribal leaders and warriors with focus on courage, hospitality \\
Love / Longing & \ar{غزل وشوق} & More personal and direct than Udhri; often expresses separation from the tribe or homeland \\
Elegy & \ar{رثاء} & Mourning for tribe members; loss of companions in the desert \\
Wisdom & \ar{حكمة} & Bedouin wisdom verses; proverbs in metered form \\
Boasting / Pride & \ar{فخر} & Tribal pride; glorification of the tribe's history and lineage \\
Camel-Driver's Call & \ar{الحداء} & A working genre: verse composed to set the rhythm of camel travel; one of the oldest Nabati purposes \\
Separation / Homeland & \ar{شعر الفراق والحنين} & Intense longing for the homeland and tribe; especially common in exile verse \\
Tribal Challenge & \ar{التحدي القبلي} & Direct tribal-vs-tribal competition; higher personal stakes than classical هجاء \\
Social Satire & \ar{هجاء شعبي} & Sharp but less courtly than classical هجاء; often humorous \\
\bottomrule
\end{tabularx}
\end{table}

% ============================================================================
\section{Regional Varieties of Dialectal Poetry}
\label{sec:regional_varieties}

\subsection{Najdi Poetry (\ar{الشعر النجدي})}

The heartland of Nabati tradition. Characterised by:
\begin{itemize}
  \item Heavy use of the Mashhub meter with its characteristic dual-rhyme system.
  \item Vocabulary drawn from Najdi dialect: phonological features include \ar{ج} → \ar{ي} in some regions (e.g., \ar{الريل} for \ar{الرجل}).
  \item Strong association with the Samri performance tradition.
  \item Famous poets: Muhsin al-Hazani, Rakan ibn Hidhlin (\ar{راكان بن حثلين}), Hamad ibn Labun (\ar{حمد بن لعبون}).
\end{itemize}

\subsection{Hijazi Poetry (\ar{الشعر الحجازي})}

The Hijaz (Mecca, Medina, Taif, Jeddah) has a distinct tradition:
\begin{itemize}
  \item The Muhawara / Qalta improvised dueling tradition is strongest here.
  \item Hijazi poets consider their tradition of poetic improvisation (\ar{المحاورة}) more refined and ancient than the Gulf version (\ar{القلطة}).
  \item Vocabulary: Hijazi dialect features including \ar{أنا} → \ar{أنا} (no guttural shift) and specific lexical items.
\end{itemize}

\subsection{Gulf Poetry (\ar{الشعر الخليجي})}

The Gulf states (Kuwait, UAE, Qatar, Bahrain) share the Nabati tradition with Najd but with local phonological and lexical features:
\begin{itemize}
  \item Gulf dialect phonology: \ar{ك} → \ar{چ} (ch) in some positions; \ar{ق} → \ar{گ} (g).
  \item The Samri tradition is also strong in Kuwait and Bahrain.
  \item Contemporary Gulf poets have taken Nabati verse to a mass broadcast audience through television programs.
\end{itemize}

% ============================================================================
\section{Famous Nabati Poets}
\label{sec:nabati_poets}

\begin{table}[H]
\centering
\caption{Landmark Nabati poets}
\label{tab:nabati_poets}
\renewcommand{\arraystretch}{1.5}
\begin{tabularx}{\textwidth}{lllX}
\toprule
\textbf{Poet} & \textbf{Arabic} & \textbf{Era} & \textbf{Contribution} \\
\midrule
Muhammad ibn Labun & \ar{محمد بن لعبون} & 18th c. & Foundational Samri poet; innovated Bahr al-Funun \\
Muhsin al-Hazani & \ar{محسن الحزاني} & 19th c. & Formalised Mashhub + dual-rhyme for Samri \\
Rakan ibn Hidhlin & \ar{راكان بن حثلين} & 19th c. & Tribal warrior poet; epic Hilali verse \\
Hamad ibn Hidhlin & \ar{حمد بن حثلين} & 19th–20th c. & Political and tribal verse \\
Muhammad ibn Shafi & \ar{محمد بن شافي} & 20th c. & Major neo-Nabati voice \\
Ghazi al-Gosaibi & \ar{غازي القصيبي} & 1940–2010 & Bridged Nabati and classical; most widely read modern Saudi poet \\
Badr ibn Abd al-Muhsin & \ar{بدر بن عبدالمحسن} & b. 1949 & Contemporary; dominant voice in Mashhub \\
\bottomrule
\end{tabularx}
\end{table}

% ============================================================================
\section{AI Implementation Notes for Nabati Poetry}
\label{sec:nabati_ai_notes}

\begin{implnote}
\textbf{Key implementation differences between Nabati and Classical:}
\begin{enumerate}
  \item \textbf{Language detection}: The system must identify the dialect (Najdi, Hijazi, Gulf) from the user's request before selecting vocabulary. Using Fusha vocabulary in a Nabati poem is a category error.
  \item \textbf{Dual-rhyme enforcement}: For the Mashhub meter, both rhyme letters must be selected and tracked separately. The classical qafiya evaluator cannot be applied to Mashhub verse.
  \item \textbf{Prosodic validation}: Nabati meters cannot be validated using the Khalilian bohour library. A separate Nabati prosody module (or LLM-based evaluation with Nabati-specific prompts) is required.
  \item \textbf{Tashkeel}: Full tashkeel is less critical in Nabati than in classical poetry, because syllable lengths are determined by dialect phonology. However, tashkeel is still useful for clarifying the intended pronunciation for prosodic scanning.
  \item \textbf{Meter selection}: Mashhub is the default for general requests. Hilali should be specified explicitly; it is perceived as more archaic and solemn.
\end{enumerate}
\end{implnote}

% ============================================================================
\clearpage
\langsep{}

% ============================================================================
%  ARABIC VERSION
% ============================================================================
\begin{otherlanguage}{arabic}

\section*{الشعر العامي والشعر النبطي}

\subsection*{الشعر النبطي: تعريفه وأصوله}

الشعر النبطي هو الشعر العامي للجزيرة العربية؛ ينظَم باللهجات المحكية — النجدية والحجازية والخليجية واليمنية — لا بالفصيح. وهو تقليد أدبي راسخ يمتد على الأقل ستة قرون، ويحظى بجمهور شعبي واسع في المملكة العربية السعودية ودول الخليج.

الشعر النبطي ليس انحدارًا عن الشعر الفصيح، بل هو نظام عروضي مستقل نما في نفس البيئة الصوتية الخليلية، غير أنه وُضع بصورة مستقلة عبر التقليد الشفهي القبلي.

\subsection*{أبرز خصائص الشعر النبطي}

\begin{enumerate}
  \item \textbf{اللغة العامية}: يُستخدم معجم اللهجة المحلية (النجدية أو الحجازية أو الخليجية)، لا الفصحى.
  \item \textbf{القافية المزدوجة}: الخاصية الأبرز في بحر المسحوب؛ لكل شطر (الصدر والعجز) قافية مستقلة، ويُشترط ألا تتقاطع القافيتان.
  \item \textbf{البحور النبطية الخاصة}: اثنا عشر بحرًا على الأقل، أشهرها: الهلالي، والمسحوب، والهجيني، والسامري، والزهيري.
  \item \textbf{الارتباط بالأداء}: الشعر النبطي في الأصل شفهي موسيقي؛ أغلبه يُغنَّى مع الربابة أو الطبول.
\end{enumerate}

\subsection*{بحر المسحوب}

بحر المسحوب هو أكثر بحور الشعر النبطي استخدامًا (أكثر من 60\% من المنظوم). اسمه مشتق من ``سحب'' النغمة مع قوس الربابة أثناء الغناء.

\textbf{تفعيلاته التقريبية:} مُسْتَفْعِلُنْ مُسْتَفْعِلُنْ فَاعِلَاتُنْ في كل شطر.

\textbf{نظام القافية المزدوجة في المسحوب:}
\begin{itemize}
  \item قافية الصدر (نهاية الشطر الأول) لها حرف روي خاص بها.
  \item قافية العجز (نهاية الشطر الثاني) لها حرف روي مختلف.
  \item يجب أن لا يتطابق حرفا الروي طوال القصيدة كلها.
  \item الخلط بين القافيتين أو توحيدهما خطأ عروضي في المسحوب.
\end{itemize}

\textbf{مثال:} يَا صَاحِبِي اللِّي يُبَالِي بِوَالِفِي --- وَالقَلْبُ مِنِّي يَتَمَنَّى لِقَاكُمْ
(قافية الصدر: الفاء --- قافية العجز: الكاف)

\subsection*{أغراض الشعر النبطي}

\begin{description}
  \item[\textbf{الشهامة والفروسية}] مدح القادة القبليين وفروسيتهم؛ نظير المديح الفصيح لكنه أكثر ارتباطًا بالحياة القبلية.
  \item[\textbf{الغزل والشوق}] الغزل في الشعر النبطي أكثر مباشرة وشخصية من الغزل الفصيح.
  \item[\textbf{الحنين والفراق}] الشوق إلى الأهل والوطن والقبيلة؛ من أكثر الأغراض التعبيرًا في الشعر النبطي.
  \item[\textbf{الحداء}] الحداء هو شعر سوق الإبل؛ شعر وظيفي يضبط إيقاع المسيرة.
  \item[\textbf{التحدي القبلي}] تنافس مباشر بين قبيلتين شعريًا؛ أكثر حدةً وشخصية من هجاء الشعر الكلاسيكي.
\end{description}

\subsection*{قواعد الذكاء الاصطناعي للشعر النبطي}

\begin{itemize}
  \item تحديد اللهجة أولًا: النجدية أم الحجازية أم الخليجية.
  \item اختيار حرفَي القافية قبل بدء التوليد (للمسحوب).
  \item لا يصلح تطبيق محقق العروض الخليلي على الشعر النبطي؛ يحتاج إلى نظام تقييم مستقل.
  \item إذا طلب المستخدم قصيدة نبطية دون تحديد البحر، يُفترض افتراضيًا المسحوب.
\end{itemize}

\end{otherlanguage}
